\setcounter{section}{-1}
\section{从一幅图开始}
本章从两道等腰直角三角形的题目出发，学习怎样把“手拉脚”“脚拉脚”转化为“手拉手”。先做题，再阅读解析；重点观察辅助点为什么取在这个位置，以及中点条件怎样参与第二次全等。题目与解析分开。



题 1 中有两个等腰直角三角形，但它们的直角顶点分别是 \(A\)、\(B\)。公共点 \(A\) 在一个三角形中是直角顶点，在另一个三角形中却是底角顶点。两个图形暂时不能绕同一个直角顶点一起旋转。

再看条件 \(BF=BD\)：\(B\) 是 \(DF\) 的中点。延长一条经过 \(B\) 的线段，往往能让 \(B\) 同时成为另一条线段的中点。这样既能出现“八字形”全等，又能补出新的等腰直角三角形。

本章的入手顺序是：\textbf{先认出两个等腰直角三角形，再检查公共点的身份；身份不一致，就用倍长或旋转补成一致。}
